\documentclass{article}
\usepackage[UTF8]{ctex}
\usepackage{amsmath,amssymb,amsfonts}
\usepackage{enumitem}
\usepackage{geometry}
\geometry{a4paper, margin=1in}
% === 错题本排版优化（不参与题目标识） ===
\setlength{\parindent}{0pt}
\setlength{\parskip}{0.35em plus 0.10em minus 0.08em}
\setlength{\emergencystretch}{3em}
\setlist{itemsep=0.25em, topsep=0.40em, parsep=0pt}
\allowdisplaybreaks[2]
% === 错题本排版优化结束 ===
\begin{document}
\section*{苏州大学2024年高等代数真题}
\section*{81．苏州大学2024年高等代数试卷}

\subsection{题目1}
一．（20 分）设 \(A=\left(\begin{array}{lll}-2 & -2 & 6 \\ -1 & -1 & 3 \\ -1 & -1 & 3\end{array}\right), C(A)=\left\{B \in \mathbb{C}^{3 \times 3} \mid A B=B A\right\}\) ．

\begin{itemize}
  \item[（1）] 求 \(A\) 的若尔当标准形．
  \item[（2）] 求复线性空间 \(C(A)\) 的维数．
\end{itemize}

\subsection{题目2}
二．（20 分）设 \(A\) 为数域 \(F\) 上的 \(n \times n\) 矩阵，\(\beta\) 为 \(F\) 上的非零列向量．证明：线性方程组 \(A X=\beta\) 有解的充分必要条件为 \(\operatorname{rank} A=\max _{1 \leq i \leq n}\left\{\operatorname{rank} A_{i}\right\}\) ，其中 \(A_{i}\) 是将 \(A\) 的第 \(i\) 列用 \(\beta\) 替换后所得矩阵。

\subsection{题目3}
三．（20分）设 \(A, B\) 为2阶方阵，\(B\) 可逆，且 \(A B=-B A\) 。证明：\(A^{2}\) 为数量矩阵。

\subsection{题目4}
四．（20分）设 \(\sigma, \tau\) 为 \(n\) 维欧氏空间 \(V\) 的线性变换，对任意的 \(\alpha, \beta \in V\) ，有

\[
(\sigma(\alpha), \tau(\beta))=(\sigma(\alpha), \beta)+(\alpha, \tau(\beta)) .
\]

证明：若 \(\mu\) 为 \(\sigma\) 的特征值，则 \(\mu \neq 1\) ，且 \(\frac{\mu}{\mu-1}\) 为 \(\tau\) 的特征值．

\subsection{题目5}
五。（20分）设 \(Q\) 为 \(n\) 阶正定矩阵，\(X\) 为实 \(n\) 维非零列向量，证明：

\begin{itemize}
  \item[（1）] \(Q+X X^{\mathrm{T}}\) 可逆；
  \item[（2）] \(0<X^{\mathrm{T}}\left(Q+X X^{\mathrm{T}}\right)^{-1} X<1\) ．
\end{itemize}

\subsection{题目6}
六．（25 分）计算题．

\begin{itemize}
  \item[（1）] 计算 \(n\) 阶行列式
\[
\left|\begin{array}{cccccc}
\cos \alpha & 1 & 0 & \cdots & 0 & 0 \\
1 & 2 \cos \alpha & 1 & \cdots & 0 & 0 \\
0 & 1 & 2 \cos \alpha & \cdots & 0 & 0 \\
\vdots & \vdots & \vdots & & \vdots & \vdots \\
0 & 0 & 0 & \cdots & 2 \cos \alpha & 1 \\
0 & 0 & 0 & \cdots & 1 & 2 \cos \alpha
\end{array}\right| .
\]
  \item[（2）] 证明：当 \(n \geq 3\) 为奇数时，\(\frac{1}{\pi} \arccos \frac{1}{\sqrt{n}}\) 是无理数。
\end{itemize}

\subsection{题目7}
七．（25 分）设 \(V\) 是数域 \(F\) 上的 \(n\) 维线性空间，\(\sigma\) 是 \(V\) 的线性变换．证明：

\begin{itemize}
  \item[（1）] 存在 \(V\) 上的线性变换 \(\tau, \tau(V)=\sigma^{-1}(0), \tau^{-1}(0)=\sigma(V)\) ．
  \item[（2）] 对任意（1）中的 \(\tau, \tau(V)=\tau^{2}(V)\) 当且仅当 \(\sigma(V)=\sigma^{2}(V)\) ．
\end{itemize}
\end{document}
